\subsection{IDENTITY ELEMENT}

DEFINITION 1. Under a law of composition \(\top\) on a set \(\mathbf{E}\) an element \(e\) of \(\mathbf{E}\) is called an identity element if, for all \(x \in \mathbf{E}\) , \(e \top x = x \top e = x\) .

There exists at most one identity element under a given law \(\top\) , for if \(e\) and \(e'\) are identity elements then \(e = e \top e' = e'\) . An identity element is permutable with every element: it is a central element.

DEFINITION 2. A magma with an identity element is called a unital magma. If E, E' are unital magmas, a homomorphism of the magma E into the magma E' which maps the identity element of E to the identity element of E' is called a unital homomorphism (or morphism) of E into E'. An associative unital magma is called a monoid.

If E, E' are monoids, a unital morphism of E into E' is called a monoid homomorphism or a monoid morphism of E into E'.

Examples. (1) In the set \(\mathbf{N}\) of natural numbers 0 is an identity element under addition and 1 is an identity element under multiplication. Each of these two laws gives \(\mathbf{N}\) a commutative monoid structure (Set Theory, III, §3, no. 3).

\begin{enumerate}
\def\labelenumi{(\arabic{enumi})}
\setcounter{enumi}{1}
\item
  In the set of subsets of a set E, \(\varnothing\) is an identity element under the law \(\cup\) and E under the law \(\cap\) . More generally, in a lattice the least element, if it exists, is identity element under the law sup; conversely, if there exists an identity element under this law it is the least element of the set. Similarly for the greatest element and the law inf.
\item
  The set \(\mathbf{N}\) has no identity element under the law \((x, y) \mapsto x^y\) . Under the law \((\mathbf{X}, \mathbf{Y}) \mapsto \mathbf{X} \circ \mathbf{Y}\) between subsets of \(\mathbf{E} \times \mathbf{E}\) the diagonal \(\Delta\) is the identity element. Under the law \((f, g) \mapsto f \circ g\) between mappings of \(\mathbf{E}\) into \(\mathbf{E}\) the identity mapping of \(\mathbf{E}\) onto \(\mathbf{E}\) is the identity element.
\item
  Let E be a magma and R an equivalence relation on E compatible with the law on E (§ 1, no. 6). If e is an identity element of E the canonical image of e in E/R is an identity element of the magma E/R.
\end{enumerate}

The identity element of a unital magma is a unital homomorphism; the composition of two unital homomorphisms is also one. For a mapping to be a unital magma isomorphism it is necessary and sufficient that it be a bijective unital homomorphism and the inverse mapping is then a unital homomorphism. Let E and E' be unital magmas and \(e'\) the identity element of \(E'\) ; the constant mapping of E into \(E'\) mapping E to \(e'\) is a unital homomorphism, called a trivial homomorphism.

The product of a family of unital magmas (resp. monoids) is a unital magma (resp. monoid).

Every quotient magma of a unital magma (resp. monoid) is a unital magma (resp. monoid).

Let E be a unital magma and e its identity element. A submagma A of E such that \(e \in A\) is called a unital submagma of E. Clearly e is the identity element of the magma A. Every intersection of unital submagmas of E is a unital submagma of E. If X is a subset of E then there exists a smallest unital submagma of E containing X; it is called the unital submagma of E generated by X; it is equal to \(\{e\}\) if X is empty. If E is a monoid, a unital submagma of E is called a submonoid of E.

If F is a magma without identity element a submagma of F may possess an identity element. For example, if F is associative and h is an idempotent element of F (I, § 1, no. 4), the set of \(h \top x \top h\) , where x runs through F, is a submagma of F with h as identity element.

If E is a magma with identity element e, a submagma A of E such that \(e \notin A\) may still possess an identity element.

DEFINITION 3. Let E be a unital magma. The identity element of E is called the composition of the empty family of elements of E.

If \((x_{\alpha})_{\alpha \in \varnothing}\) is the empty family of elements of E its composition \(e\) is also denoted by \(\bigcap_{\alpha \in \varnothing} x_{\alpha}\) . For example, we write

\[
\bigcap_ {q \leqslant i \leqslant p} x _ {i} = e
\]

when \(p < q\) ( \(p, q \in \mathbf{N}\) ). Similarly we write \(\top^0 x = e\) for arbitrary \(x\) . With these definitions Theorems 1 and 3 of § 1 remain true if the hypothesis that the sets \(\mathbf{A}\) and \(\mathbf{B}_i\) are non-empty is suppressed. Similarly the formulae \(\top^{m+n} x = \left( \frac{m}{\top} x \right)^{\top} \left( \frac{n}{\top} x \right)\) and \(\top^{mn} x = \frac{m}{\top} \left( \frac{n}{\top} x \right)\) are then true for \(m \geqslant 0, n \geqslant 0\) .

Let E be a unital magma whose law is denoted by \(\top\) and \(e\) its identity element. The support of a family \((x_{i})_{i\in I}\) of elements of E is the set of indices \(i\in I\) such that \(x_{i}\neq e\) . Let \((x_{i})_{i\in I}\) be a family of elements of E with finite support. We shall define the composition \(\bigcap_{i\in I}x_{i}\) in the two following cases:

\begin{enumerate}
\def\labelenumi{(\alph{enumi})}
\item
  the set I is totally ordered;
\item
  E is associative and the \(x_{i}\) are pairwise permutable.
\end{enumerate}

In these two cases let S be the support of the family \((x_{i})\) . If J is a finite subset of I containing S, then \(\bigcap_{i\in J}x_{i} = \bigcap_{i\in S}x_{i}\) , as is seen by induction on the number of elements in J, applying Theorem 1 of §1 in case (a) and Theorem 3 of §1 in case (b). Let \(\bigcap_{i\in I}x_{i}\) denote the common value of the compositions \(\bigcap_{i\in J}x_{i}\) for all finite subsets of I containing S. When I is the interval \(\{p,\rightarrow\{\text{of }N\}\) , we also write \(\bigcap_{i=p}^{\infty}x_{i}\) .

With these definitions and notation, Theorems 1 and 3 of § 1 and the remarks following Theorem 3 extend to families with finite support.

The identity element of a law written additively is usually denoted by 0 and called zero or the null element (or sometimes the origin). Under a law written multiplicatively it is usually denoted by 1 and called the unit element (or unit).

